\documentclass[11pt,a4paper]{article}
\usepackage[utf8]{inputenc}
\usepackage[T1]{fontenc}
\usepackage[margin=2.5cm]{geometry}
\usepackage{amsmath}
\usepackage{amssymb}
\usepackage{xcolor}
\usepackage{soul}
\usepackage{hyperref}
\usepackage{enumitem}
\usepackage{parskip}

\definecolor{revblue}{RGB}{0,70,140}
\definecolor{revgreen}{RGB}{0,120,60}
\definecolor{revgray}{RGB}{90,90,90}
\definecolor{revred}{RGB}{180,40,40}

\hypersetup{colorlinks=true,linkcolor=revblue,urlcolor=revblue,citecolor=revblue}

% Custom command for reviewer quote box
\newenvironment{revquote}{%
  \begin{quote}\small\itshape\color{revgray}
}{%
  \end{quote}
}

% Custom command for our response
\newcommand{\response}[1]{\par\noindent\textbf{\color{revgreen}Authors' response:} #1\par}

% Custom command for change reference
\newcommand{\changes}[1]{\par\noindent\textbf{\color{revblue}Changes in manuscript:} #1\par}

% Status markers
\newcommand{\done}{\color{revgreen}\textbf{[DONE]}}
\newcommand{\inprogress}{\color{revblue}\textbf{[IN PROGRESS]}}
\newcommand{\pending}{\color{revred}\textbf{[PENDING]}}

\title{\textbf{Point-by-Point Response to Reviewers} \\
\large Manuscript: ``Same Prompt, Different Answer: Hidden Non-Determinism in LLM APIs Undermines Scientific Reproducibility'' \\
\large Submitted to \emph{Nature Communications}}

\author{Lucas Rover, Hugo Valadares Siqueira, Eduardo Tadeu Bacalhau, \\
Anibal Tavares de Azevedo, Yara de Souza Tadano \\
Federal University of Technology --- Paraná (UTFPR)}

\date{Revision submitted: TBD --- Original submission: 2026-03-12}

\begin{document}

\maketitle

\section*{Cover note to the Editor}

Dear Dr.\ Bigorajski,

We thank you and the three reviewers for the careful and constructive engagement with our manuscript. Your decision to invite a major revision, together with the reviewers' clearly articulated concerns, has substantially improved the work. We have addressed every point raised, with the following major additions:

\begin{enumerate}[leftmargin=*]
  \item \textbf{Conceptual reframe to deployment stacks.} Following Reviewer~1's central observation, we have reorganised the manuscript around the concept of a \emph{deployment stack} --- the tuple (model weights, provider, serving infrastructure, API layer) --- as the unit of analysis. This reframe is now explicit in the title (subtitle), abstract, introduction, methods, tables, figure captions, and discussion. The Together AI quasi-isolation result, which previously appeared as a single empirical observation, is now positioned as the mechanistic core of the paper: it demonstrates directly that the stack, not the abstract model, is the carrier of variation.
  \item \textbf{Expansion to coding, math, and out-of-domain tasks.} Following the editor's specific request and Reviewers~1 and~3, we ran new controlled experiments on HumanEval (30 problems), GSM8K (30 problems), and 10 PubMed PM\textsubscript{2.5} abstracts. Across these new domains, the local--API gap persists and widens precisely where small differences matter: on GSM8K math reasoning, Claude Sonnet~4.5 reaches EMR~=~0.063 [0.03, 0.10]; on PubMed PM\textsubscript{2.5} extraction, Claude reaches EMR~=~0.010 [0.00, 0.03]; on HumanEval code generation, Claude reaches EMR~=~0.39 [0.27, 0.52] while local stacks maintain EMR~$\geq$~0.92. The pattern is domain-transferable, not corpus-specific.
  \item \textbf{Field-level reproducibility analysis.} Following Reviewer~1, we now compute every reproducibility metric (EMR, NED, ROUGE-L, BERTScore F1) per JSON output field. The results reveal that BERTScore is structurally saturated and unable to discriminate between conclusion-relevant and metadata fields, while EMR exposes a paired Cohen's~$d=+1.41$ gap between conclusion-relevant fields ($\overline{\text{EMR}}=0.455$) and metadata fields ($\overline{\text{EMR}}=0.684$). This resolves the apparent tension Reviewer~1 identified between ``textual not semantic'' and ``substantive content'' --- the resolution is that BERTScore alone is the wrong instrument; the three-level framework we propose is required to detect substantive payload variation.
  \item \textbf{Multi-turn extension to gpt-4o and DeepSeek.} Following Reviewer~3 (item~3), we extended the three-turn refinement protocol to OpenAI gpt-4o and DeepSeek Chat (50 runs each). The near-zero reproducibility under multi-turn previously reported for Claude (0.04) and Gemini (0.01) is confirmed for gpt-4o (EMR~=~0.09 [0.02, 0.16]) and partially for deepseek-chat (EMR~=~0.35 [0.13, 0.60], down from 0.84 single-turn --- a four-fold drop across dialogue turns). The pattern is a general property of cloud-served interactive pipelines, not a Claude/Gemini-specific anomaly.
  \item \textbf{Applied evidence-synthesis impact and validation.} Following Reviewer~1, we promote the PM\textsubscript{2.5}/respiratory health analysis to a dedicated Results subsection (``Applied impact in evidence synthesis''). To preserve the contributions of our companion paper currently submitted (Rover \& Tadano, 2026, \emph{Research Synthesis Methods}), the present manuscript reports only the aggregate propagation finding and cites the companion paper for the detailed pairwise reproducibility, silver-standard validation, and meta-analytic per-run analyses. For a self-contained in-paper triangulation of the divergence claim (R3.6), we ran a two-judge LLM-as-judge protocol (Claude Opus~4.7 and gpt-4o) on a \emph{new} 30-case sample (distinct from the 23 effects analysed in the companion paper) with three pre-registered criteria. \textbf{Both judges classify the majority of divergences as truly contradictory} (Claude 22/30 = 73\% [56--86\%]; gpt-4o 27/30 = 90\% [74--97\%]). Inter-judge $\kappa$~=~0.29 (fair). The majority finding is therefore robust to the choice of LLM judge and BERTScore F1~$>$~0.97 on the same outputs cannot detect this substantive payload divergence---confirming the R1.5 finding.
\end{enumerate}

The revised manuscript with all changes marked is provided alongside this response, together with the revised supplementary information. We have followed Nature Communications' formatting and reporting requirements, including the Code/Software submission checklist, the Machine Learning checklist, the Reporting Summary, the Data Availability and Code Availability sections, and ORCID linking for all corresponding authors. All experimental records are deposited at the project repository (\url{https://github.com/Roverlucas/genai-reproducibility-protocol}, archived to Zenodo on acceptance) and on figshare (DOI: \url{https://doi.org/10.6084/m9.figshare.31653373}).

We hope the revisions meet the high standard you and the reviewers have set. Detailed responses to every comment follow below; reviewer comments are reproduced verbatim in italic grey.

\bigskip

Yours sincerely,

\medskip

\textbf{Lucas Rover} \\
on behalf of all authors \\
\href{mailto:lucasrover@alunos.utfpr.edu.br}{lucasrover@alunos.utfpr.edu.br}

\newpage
%==============================================================
\section*{Response to Reviewer 1}
%==============================================================

We thank Reviewer~1 for the comprehensive and unusually constructive review. The two strengths the reviewer identified --- the three-level reproducibility framework and the temperature-paradox observation --- are central to the contribution we wished to make, and we are grateful that they were recognised. The fifteen substantive issues the reviewer raised have led to substantial improvements; each is addressed individually below.

%--------------------------------------------------------------
\subsection*{R1.1 --- Framing: APIs and infrastructure stacks, not models}
%--------------------------------------------------------------

\begin{revquote}
While the paper frames the testing as testing of different models, what is being tested is actually APIs and associated infrastructure stacks. Specifically, the same model served by a different API may have different degrees of non-determinism. Could this framing be corrected?
\end{revquote}

\response{We agree entirely. This is the central conceptual observation of the review and we have restructured the manuscript accordingly. We now define a \emph{deployment stack} as the tuple (model weights, provider, serving infrastructure, API layer) and treat it as the unit of analysis throughout. Each component of the tuple can independently introduce variation, and the same model weights served by two different stacks constitute two distinct experimental objects for reproducibility purposes. The Together AI quasi-isolation result --- in which the same LLaMA~3 8B weights served via two different stacks produced markedly different reproducibility profiles --- is now positioned as the direct empirical evidence that the stack, not the model, carries the variation.} \done

\changes{
\begin{itemize}[leftmargin=*]
  \item \textbf{Abstract} (lines 62--67): inserted ``deployment stacks'' tuple definition; substituted ``API-served models'' $\rightarrow$ ``API-served stacks''; clarified the quasi-isolation as same weights served by two different stacks.
  \item \textbf{Introduction} (lines 99--101, 106): new $\sim$80-word paragraph anchoring the deployment-stack concept; rewritten third contribution to emphasise ``same weights served via two stacks''.
  \item \textbf{Results §1 heading and lead} (lines 115, 117): renamed and updated terminology consistently.
  \item \textbf{Table~1} (lines 240--262): caption rewritten; column header ``Model / Deployment'' $\rightarrow$ ``Deployment stack (weights) / Class''; cross-reference to Methods §4.2 added.
  \item \textbf{Figures 1, 2, 5 captions} (lines 230--235, 271--276, 315--318): ``model deployments'' $\rightarrow$ ``deployment stacks''; ``local models / API models'' $\rightarrow$ ``local stacks / API stacks''.
  \item \textbf{Methods --- new subsection ``Unit of analysis: deployment stacks''} (lines 407--415): $\sim$155-word subsection inserted before ``Models and infrastructure'' (now renamed ``Deployment stacks evaluated''), defining the tuple and its operational interpretation throughout the manuscript.
  \item \textbf{Methods --- ``Deployment stacks evaluated''} (line 417 onwards): renamed from ``Models and infrastructure''; entries restructured as full tuples (e.g., L-LLaMA3 = (LLaMA~3 8B Q4 / local Ollama / single-GPU Apple~M4 / \texttt{/api/generate} endpoint)).
  \item \textbf{Discussion} (lines 341--344): new $\sim$110-word paragraph after the Together AI quasi-isolation discussion, explicitly stating ``the deployment stack, not the abstract model, is the carrier of variation in LLM-based research''.
\end{itemize}
}

%--------------------------------------------------------------
\subsection*{R1.2 --- Comparability across providers with different parameter exposure}
%--------------------------------------------------------------

\begin{revquote}
The authors write that ``all providers expose the same user-facing `deterministic' parameters (temperature zero, fixed seed where supported)'' yet later note that the settings each specific API offers vary. How does this affect comparability?
\end{revquote}

\response{The reviewer correctly identifies that providers expose overlapping but not identical parameter surfaces. The deterministic configuration we set is the maximum each provider exposes for that purpose. We now make this explicit. Specifically: temperature zero is universally exposed; the seed parameter is supported by OpenAI (advisory) and Gemini, absent from Anthropic, and not exposed by DeepSeek's OpenAI-compatible Chat Completions interface or by Perplexity Sonar. Our \texttt{seed\_status} Run-Card field already records this asymmetry per stack (\texttt{sent}, \texttt{logged-only-not-sent}, or \texttt{not-supported}). What is comparable across stacks is therefore not parameter \emph{equality} but the operational claim that we used the maximum determinism configuration each provider documents, with the precise per-stack parameter set explicit in Supplementary §S4 and the Run Card.} \done

\changes{Methods §``Deployment stacks evaluated'' (line 417 onwards) now lists per-stack parameter exposure explicitly. Supplementary §S4 already documented the per-stack payloads; we have added a small introductory paragraph clarifying that comparability is at the level of ``maximum determinism configuration documented by each provider'' rather than parameter-by-parameter equality.}

%--------------------------------------------------------------
\subsection*{R1.3 --- Perplexity Sonar's search-augmented architecture (line 103)}
%--------------------------------------------------------------

\begin{revquote}
Line 103 states that ``Perplexity Sonar's particularly low reproducibility reflects its search-augmented architecture, where real-time web retrieval introduces an additional source of variation'' --- it is unclear whether this claim is directly supported by the paper's experiments. If not, the statement should either be supported by appropriate citation or explicitly framed as informed speculation.
\end{revquote}

\response{The reviewer is correct that this claim is not directly tested by an isolation experiment in our study. We have reframed it as an informed hypothesis with citations to Perplexity's public Sonar documentation and to the broader literature on retrieval-augmented generation variability. The reformulated text now reads: ``\ldots\ is consistent with its publicly documented search-augmented architecture [perplexitysonar2024]\ldots\ A plausible mechanism is that variability in the retrieved document set compounds with model-internal non-determinism: the foundational RAG framework [lewis2020rag]\ldots\ [shi2024ragvariability]\ldots\ We did not perform a controlled retrieval-isolation experiment for Sonar (the index is not user-accessible), so this mechanism is offered as an informed hypothesis rather than a demonstrated cause.''} \done

\changes{Results §``Non-determinism varies substantially across providers'' (around line 158): Perplexity passage rewritten. Three new \texttt{\textbackslash bibitem} entries appended to the inline bibliography: \texttt{lewis2020rag}, \texttt{perplexitysonar2024}, \texttt{shi2024ragvariability}.}

%--------------------------------------------------------------
\subsection*{R1.4 --- Cloud deployment vs.\ production serving infrastructure}
%--------------------------------------------------------------

\begin{revquote}
Line 110 states that ``A natural hypothesis is that cloud deployment itself (network latency, load balancing, shared hardware) causes the non-determinism'', while line 120 states ``The variability observed in GPT-4, Claude, and Gemini is instead consistent with the complexity of their production serving infrastructure''. The distinction between ``cloud deployment itself'' and ``production serving infrastructure'' is currently fuzzy and would benefit from clearer distinction. Listing facets of cloud deployment and known sources of non-determinism is helpful. However, it is not clear that it is comprehensive enough to support the strong conclusions being drawn. What additional factors might be at play in both cloud infrastructure and inference infrastructure?
\end{revquote}

\response{We thank the reviewer for pressing this point; the distinction is indeed underspecified. We have added a new opening paragraph to Methods §``Sources of non-determinism in distributed inference'' that explicitly distinguishes the two factor classes: ``\emph{Cloud deployment factors} are those intrinsic to running a model on hardware accessed over a network: round-trip latency, request routing across geographically distributed endpoints, shared physical hardware, and multi-tenancy isolation. \emph{Production serving infrastructure factors} are specific design choices in the LLM inference stack: tensor parallelism with non-deterministic all-reduce, FlashAttention kernel non-determinism, dynamic batching with continuous request scheduling, speculative decoding fast-paths, mixed-precision BF16/FP16 accumulation, and prefix-caching state across requests.'' The Together~AI quasi-isolation probe controls for the first class (the request still traverses the network, hits a multi-tenant cluster) while differing in serving-infrastructure complexity from the local single-GPU baseline; that this configuration achieves near-local reproducibility is what licenses the inference that the gap originates in serving infrastructure rather than cloud deployment per se. This is now stated explicitly in the new paragraph.} \done

\changes{Methods §``Sources of non-determinism in distributed inference'' (around line 553): new $\sim$120-word opening paragraph explicitly disambiguating cloud-deployment vs production-serving-infrastructure factors, with the Together~AI result invoked to license the inferential bridge. See also the new mechanism mapping table (R3.5).}

%--------------------------------------------------------------
\subsection*{R1.5 --- Field-level analysis: textual vs substantive divergence}
%--------------------------------------------------------------

\begin{revquote}
Line 172 states that ``The key result field diverges in 67\% of groups, and method in 57\%''. However, it is not clear whether differences in these fields implies substantive changes in content. Given the reported high semantic similarity and low text edit distance, these differences might reflect minor wording variations than meaningful semantic changes. As currently written, the manuscript appears to conflate opposing claims ``difference is textual not semantic'' and ``difference affects the substantive content'' that need to be resolved.

[\ldots] It would strengthen the argument if the authors computed all metrics on these specific fields to test whether semantic similarity is indeed lower there.
\end{revquote}

\response{We thank the reviewer for identifying this tension and prompting the per-field analysis. We have computed EMR, NED, ROUGE-L, and BERTScore F1 separately for each of the five extracted fields (objective, method, key\_result, model\_or\_system, benchmark) across all stacks. The result resolves the apparent contradiction in a way we did not anticipate and that, on reflection, strengthens the paper:

\begin{itemize}[leftmargin=*]
  \item Across the API stacks where outputs actually diverge, the conclusion-relevant fields (objective, method, key\_result) average $\overline{\text{EMR}}=0.455$, while the metadata fields (model\_or\_system, benchmark) average $\overline{\text{EMR}}=0.684$. This is a paired Cohen's $d=+1.41$ effect --- a very large gap.
  \item However, BERTScore F1 is \emph{saturated} across both groups: conclusion-relevant fields average 0.978 and metadata fields 0.979, paired $d=-0.10$ (effectively null).
\end{itemize}

The interpretation is therefore: \emph{BERTScore is structurally unable to discriminate substantive payload divergence from cosmetic textual variation} when one operates on extracted-JSON fields of moderate length. EMR and NED can. The single most striking single-stack illustration is Gemini~2.5~Pro RAG on \texttt{key\_result}, where EMR$=0.10$ (the field changes 90\,\% of the time) while BERTScore F1 remains 0.969. We therefore now argue that the apparent tension the reviewer identified is itself evidence that the three-level reproducibility framework (bitwise $\to$ structural $\to$ semantic) we propose is required: any single metric, including a state-of-the-art semantic similarity metric, would mislead.} \done

\changes{
\begin{itemize}[leftmargin=*]
  \item New \textbf{Extended Data Table~7-bis} (\texttt{analysis/tables/table\_per\_field\_metrics.tex}) reporting per-stack $\times$ per-field EMR, NED, ROUGE-L F1, and BERTScore F1 with 10\,000-resample percentile bootstrap CIs.
  \item New \textbf{Extended Data Figure} (\texttt{analysis/figures/per\_field\_radar.pdf}) showing per-field BERTScore F1 by stack with a vertical separator between conclusion-relevant and metadata fields, illustrating BERTScore saturation.
  \item Revised \textbf{Results §``Outputs diverge textually but not semantically''} now reports the paired Cohen's $d=+1.41$ EMR gap and the BERTScore null result, and reframes the argument explicitly: BERTScore alone is the wrong instrument for substantive-payload reproducibility, and the three-level framework is necessary precisely because of this.
  \item Discussion now contains the Gemini~2.5~Pro RAG single-stack illustrative case (EMR$=0.10$, BERTScore F1$=0.969$).
\end{itemize}
}

%--------------------------------------------------------------
\subsection*{R1.6 --- Propagation: textual differences affect automated processing}
%--------------------------------------------------------------

\begin{revquote}
``Where outputs are parsed programmatically rather than read by humans, even minor lexical differences can propagate''. Paired with the next paragraph, the point seems to be that these very minor, mostly textual differences can nonetheless affect automated processing of the outputs. This is a key part of the argument, as it prevents the findings from being dismissed as trivial or irrelevant given the high semantic similarity of the results. Clarifying this connection would strengthen the claim that ``non-determinism is therefore not a technical curiosity but a property that interacts with how LLMs are deployed as measurement instruments''.
\end{revquote}

\response{We agree this is the central argumentative bridge that prevents the findings from being dismissed as cosmetic. We have rewritten the paragraph to make the propagation logic explicit: (i)~when outputs feed automated parsers (regex extractors, JSON schema validators, downstream classifiers), even single-character differences in numeric values, dates, or named entities cause divergent extracted records; (ii)~when downstream pipelines aggregate these records into statistics (means, counts, effect sizes), each individual lexical divergence becomes a distinct data-point; (iii)~the per-field EMR drop on \texttt{key\_result} (R1.5) is exactly where this propagation is most consequential, because the \texttt{key\_result} field is what evidence-synthesis pipelines aggregate. The companion PM\textsubscript{2.5} analysis (now expanded) shows this directly: 23 distinct effect estimates appear or disappear from the evidence base depending solely on which run of the LLM is used.} \done

\changes{New Results subsection ``Applied impact in evidence synthesis: an out-of-AI/ML probe'' inserted before the Discussion (manuscript lines 246--264). It explicitly chains the propagation logic from the per-field analysis (R1.5) to the 10-abstract PubMed PM\textsubscript{2.5} extraction probe and, via citation of the companion paper, to the 23 study-level effect estimates that appear or disappear depending on the LLM run.}

%--------------------------------------------------------------
\subsection*{R1.7 --- Figure 1 caption inconsistency (8 vs 5 vs 10 deployments)}
%--------------------------------------------------------------

\begin{revquote}
Fig~1. The caption refers to 8 model deployments but the figure appears to show 5 or 10 depending on how ``deployment'' is defined (i.e the number of models or number of experiments). Additionally, the results for deepseek and perplexity are not shown. More broadly, why do most figures report results for only two API models, and why do the specific models vary across figures?
\end{revquote}

\response{We thank the reviewer for catching this; the caption was inconsistent with the rendered figure. We have rewritten Figure~1's caption to enumerate all eight stacks shown (3 local + Together~AI + 4 API) and to explicitly note that API-Gemini appears only in Figure~2 because it has Tasks~3--4 coverage only. Coverage variation across figures was driven by data availability per task --- e.g., DeepSeek and Perplexity were evaluated under C1 only on Tasks 1--2 due to API-quota constraints documented in Supplementary §S10. We now state this coverage matrix explicitly in the captions of every figure that reports a subset.} \done

\changes{Figure~1 caption rewritten (around lines 230--235) enumerating all eight stacks shown and explicitly explaining the Gemini absence; cross-reference to Supplementary §S10 added. Figs~2 and~5 captions also reviewed for stack-list accuracy.}

%--------------------------------------------------------------
\subsection*{R1.8 --- Code generation and mathematical reasoning experiments}
%--------------------------------------------------------------

\begin{revquote}
As one of the limitations, the authors specify that code generation and mathematical reasoning are excluded from the set of tasks. However, running experiments on code generation and mathematical reasoning could greatly strengthen the results because these are domains where minor differences in text actually can completely make or break the functionality of the output. The authors could sample tasks from coding benchmarks as a quick and easy way to test this.
\end{revquote}

\response{We thank the reviewer for the suggestion, which the editor also identified as a precondition for resubmission. We ran new controlled experiments under our standard C1 greedy-decoding protocol ($t{=}0$, fixed seed) across all eight deployment stacks (3 local + Together AI quasi-isolation + 4 API):

\textbf{HumanEval (code generation, 30 problems, 5 reps per problem).} EMR results:
\begin{itemize}[leftmargin=*]
  \item Local stacks (Ollama on M4) + Together AI: EMR $\geq$ 0.92, with Gemma~2~9B, Together AI, and Gemini~2.5~Pro reaching the perfect 1.000.
  \item API stacks: deepseek-chat = 0.84, gpt-4o = 0.84, \textbf{Claude Sonnet~4.5 = 0.39 [0.27, 0.52]}.
\end{itemize}

\textbf{GSM8K (math reasoning, 30 problems, 5 reps per problem).} EMR results:
\begin{itemize}[leftmargin=*]
  \item Local + Together + Gemini: EMR $\geq$ 0.84, with three stacks reaching 1.000.
  \item API stacks fall sharply: deepseek-chat = 0.37, gpt-4o = 0.27, and \textbf{Claude Sonnet~4.5 = 0.063 [0.03, 0.10]}.
\end{itemize}

The local-to-API gap is therefore not a summarisation-task artefact. It widens precisely on tasks where small textual differences alter functionality (HumanEval --- a single-character change can break unit tests) or numerical correctness (GSM8K --- one digit changes the final answer). The Claude Sonnet 4.5 EMR of 0.063 on GSM8K is particularly striking: under documented-deterministic settings, character-identical reproduction occurs on roughly one in 16 paired runs of the same math problem.} \done

\changes{New Results subsection ``Coding and math reasoning'' with EMR table per stack (above), referencing the new Extended Data tables \texttt{table\_t1\_t4\_t14.tex} (auto-generated from \texttt{analysis/revision/emr\_per\_stack\_per\_task.json}). Limitations section updated to remove the original ``code and math excluded'' caveat.}

%--------------------------------------------------------------
\subsection*{R1.9 --- Greater emphasis on the applied PM2.5 experiment}
%--------------------------------------------------------------

\begin{revquote}
The applied experiment (line 270--277) is crucial for the paper's central claim, in particular the argument about the downstream effect textual non determinism has on automated processing and ultimately affect scientific work that depends on these methods. Could the authors clarify why this part is not given greater emphasis within the overall study given its importance for demonstrating downstream impact?
\end{revquote}

\response{We agree the applied PM\textsubscript{2.5}/respiratory health analysis warrants prominent placement and we have promoted it to a dedicated Results subsection, ``Applied impact in evidence synthesis: an out-of-AI/ML probe'' (manuscript lines 246--264). The section reports the 10-abstract PubMed PM\textsubscript{2.5} extraction probe (Claude EMR~=~0.010, mirroring the 0.020 observed on AI/ML summarisation, see R3.4) and the in-paper two-judge LLM-judge triangulation of 30 newly sampled cases (see R3.6, verdicts Claude 22/3/5; gpt-4o 27/3/0). To avoid duplicate publication with our companion paper currently submitted to \emph{Research Synthesis Methods} (Rover \& Tadano, 2026), the present manuscript reports only aggregate statistics, the propagation finding (up to 23 study-level effect estimates appear or disappear depending solely on which LLM run is used), and the triangulation outcome, with full per-abstract data and analysis pipeline residing in the companion paper, which is cited.} \done

\changes{New Results subsection ``Applied impact in evidence synthesis: an out-of-AI/ML probe'' (manuscript lines 246--264); new Methods Table~1 (label \texttt{tab:revision\_emr}) reporting per-stack EMR on the 10-abstract probe; in-paper LLM-judge triangulation verdicts (5/3/2) reported in-line. Companion paper cited but its detailed pairwise, silver-standard, and meta-analytic tables are not reproduced here.}

%--------------------------------------------------------------
\subsection*{R1.10 --- Protocol scope: client-side vs provider-side}
%--------------------------------------------------------------

\begin{revquote}
Protocol design described in line 310 ``Run Cards capture the complete execution context of a single generative AI run'' appears to primarily be relevant only for open-source settings. Could the authors clarify whether they envision proprietary API providers adopt this, for example, documenting details such as the weights hash and environment fingerprint to enable provenance tracking? In practice, the weights themselves and full documentation of production serving infrastructure are not generally made public. More fundamentally, if the main source of non determinism is how production infrastructure itself functions, it is unclear how documenting that infrastructure would resolve reproducibility issues. The paper already demonstrates that using identical infrastructure does not guarantee identical outputs. It would thus be helpful to clarify what exactly documenting the environment changes.
\end{revquote}

\response{We thank the reviewer for surfacing an ambiguity in our framing. The provenance protocol is, by design, a \emph{client-side} instrument. We have inserted a new paragraph titled ``Scope: client-side observability'' at the start of Methods §``Protocol design'': ``It documents what is observable from the client side: the request payload, the resolved model identifier returned by the API, response headers, the API response identifier, latency, and the input/output text together with their cryptographic hashes. It does not require provider-side weight hashes\ldots\ For proprietary APIs in particular, the protocol creates the audit trail that providers' opacity would otherwise make impossible. The protocol is therefore not restricted to open-source stacks; it is the minimum machinery that closes the audit loop on any deployment stack.'' The protocol's operational value is twofold: (i)~differential diagnosis (when the same prompt produces different outputs at the same snapshot, the logged hashes establish the divergence is attributable to the stack rather than client-side drift); and (ii)~provider-agnostic minimum standard (deploying the protocol creates the audit trail that opacity otherwise prevents). Whether a provider extends the protocol with weight hashes or deterministic-execution attestation is a separate question of provider commitment; the protocol functions in either case.} \done

\changes{Methods §``Protocol design'' (around line 440): new $\sim$150-word opening paragraph titled ``Scope: client-side observability'' inserted, making the client-side framing explicit and pre-empting the open-source restriction reading.}

%--------------------------------------------------------------
\subsection*{R1.11 --- Anthropic seed clarification (line 363)}
%--------------------------------------------------------------

\begin{revquote}
Line 363 states ``For API models, the seed parameter is advisory (OpenAI[12]), absent (Anthropic), or empirically insufficient (Gemini)''. For Anthropic, do[es] this imply that even under the C1 condition, the seed may not in fact be fixed, and could therefore be a source of non-determinism? Or is the point that the authors used an Anthropic API that Anthropic claimed would be deterministic, but observed that it still produced non-deterministic outputs in reality?
\end{revquote}

\response{We thank the reviewer for catching this ambiguity. The Anthropic Messages API does not accept a seed parameter at all, so no seed value is transmitted with any Claude request. Our Run Card nevertheless records \texttt{seed:\ 42} together with \texttt{seed\_status:\ "logged-only-not-sent-to-api"} for every Claude run, so that downstream audit tools can distinguish a seed that was requested by the experimenter from one that was honoured by the API (Supplementary §S4). Claude is consequently evaluated under \texttt{temperature\ =\ 0}, which Anthropic documents as producing the most-probable token. The non-determinism we observe for the Claude stack is therefore attributable to the production serving infrastructure rather than to seed variation, because seed variation is not even a possibility on the Anthropic surface. The clarification is now in Methods §``Experimental conditions''.} \done

\changes{Methods §``Experimental conditions'' (around line 526): paragraph appended with the Anthropic-seed clarification and \texttt{seed\_status} explanation. Cross-reference to Supplementary §S4 added.}

%--------------------------------------------------------------
\subsection*{R1.12 --- Concrete examples of observed differences}
%--------------------------------------------------------------

\begin{revquote}
It would have been helpful if the authors included concrete examples of the observed differences to supplement the metrics.
\end{revquote}

\response{We agree, and we have added a new Box (Box~1) with three concrete side-by-side examples drawn directly from real Run Cards in the public repository: (i)~GPT-4 extraction abstract\_001 \texttt{benchmark} field --- a comma-vs-``and'' textual difference (cosmetic); (ii)~Claude RAG abstract\_001 (reps~0 vs~1) \texttt{key\_result} field --- ``parallelizability + 8 GPUs'' vs ``training time + constituency parsing generalisation'' (substantive content change); (iii)~same group, \texttt{method} field --- fully reformulated wording (substantive). Diff highlights use coloured bold to mark the divergent tokens. The Box is referenced in the main text after Fig.~5 and is reproducible from the open repository (Run IDs cited in the figure caption).} \done

\changes{New Box~1 (implemented as a \texttt{figure*} float wrapping \texttt{\textbackslash fbox\{\textbackslash begin\{minipage\}\dots\}}) inserted after Fig.~5 (\texttt{fig:overview}) so existing fig refs are not renumbered. In-text reference: ``(Box~1, shown after Fig.~\textbackslash ref\{fig:overview\})''. Three examples drawn from real Run Cards.}

%--------------------------------------------------------------
\subsection*{R1.13 --- W3C PROV definition (line 319)}
%--------------------------------------------------------------

\begin{revquote}
Line 319 could benefit from an additional sentence defining what W3C PROV is and why it's being used.
\end{revquote}

\response{Done. We have added a definition on first use of ``W3C PROV'' in the introduction (around line~107): ``the W3C Provenance (PROV) Data Model [w3cprov2013] --- a W3C Recommendation that defines an interoperable standard for representing the entities, activities, and agents involved in producing a piece of data\ldots''.} \done

\changes{Introduction (around line 107): definition sentence added on first use of ``W3C PROV''. Citation key \texttt{w3cprov2013} reused (already in the inline bibliography).}

%--------------------------------------------------------------
\subsection*{R1.14 --- Reporting Summary ``deployment mode''}
%--------------------------------------------------------------

\begin{revquote}
The Reporting summary, in ``A description of all covariates tested'' list ``deployment mode''. Could the authors clarify whether this refers to the decoding mode and where they vary the prompt format?
\end{revquote}

\response{Thank you for catching this. ``Deployment mode'' refers to the local-vs-API axis (whether a stack is served locally on a single GPU via Ollama or via a cloud provider API) and is a stack-level covariate fixed for the lifetime of each experimental group, not a within-run varying parameter. The revised Reporting Summary entry now reads explicitly: ``deployment mode refers to whether a model is served locally (single-GPU Ollama on Apple~M4) or via a remote cloud API\ldots\ a stack-level covariate fixed for the lifetime of each experimental group, not a per-run varying parameter.'' This distinguishes it from decoding mode (greedy at $t=0$, fixed across deterministic conditions C1--C2) and from prompt format (controlled in Supplementary §S7 with the chat-format control experiment).} \done

\changes{\texttt{article/reporting\_summary\_filled.tex} line~56 (covariates row): revised wording inserted; cross-reference to Supplementary §S7 added.}

%--------------------------------------------------------------
\subsection*{R1.15 --- API documentation links in Supplementary §S4}
%--------------------------------------------------------------

\begin{revquote}
S4 in the Supplementary Information could benefit from adding a link to each API's official documentation. This would further disambiguate and clarify exactly what was used, and ensure clarity and reproducibility in case API names and descriptions change over time.
\end{revquote}

\response{Done. Supplementary §S4 now includes a new Table~S2-bis with hyperlinks to each provider's official API reference: OpenAI (\url{https://platform.openai.com/docs/api-reference/chat}), Anthropic (\url{https://docs.anthropic.com/en/api/messages}), Google Gemini (\url{https://ai.google.dev/api/rest}), DeepSeek (\url{https://api-docs.deepseek.com}), Perplexity (\url{https://docs.perplexity.ai}), Together~AI (\url{https://docs.together.ai/reference/chat-completions}), and Ollama (the local stack). Access date is recorded.} \done

\changes{Supplementary §S4: new \texttt{Table~\textbackslash ref\{tab:apidocs\}} appended with seven hyperlinks (six providers requested by the reviewer plus Ollama for completeness).}

\newpage
%==============================================================
\section*{Response to Reviewer 2}
%==============================================================

\begin{revquote}
I co-reviewed this manuscript with one of the reviewers who provided the listed reports. This is part of the Nature Communications initiative to facilitate training in peer review and to provide appropriate recognition for Early Career Researchers who co-review manuscripts.
\end{revquote}

\response{We thank Reviewer 2 for participating in the co-review and welcome the Early Career Researcher initiative. We note that no independent comments were provided.}

\newpage
%==============================================================
\section*{Response to Reviewer 3}
%==============================================================

We thank Reviewer 3 for the careful and constructive review and for explicitly identifying the quasi-isolation probe as the primary novelty of the work. The six issues raised are addressed individually below.

%--------------------------------------------------------------
\subsection*{R3.1 --- More creative/open-ended tasks}
%--------------------------------------------------------------

\begin{revquote}
The primary conclusions are drawn from scientific summarization and JSON extraction. More creative or open-ended tasks (code generation, math, creative writing) should be included to verify the extent of ``hidden'' variation in those critical domains.
\end{revquote}

\response{We thank the reviewer for the suggestion, which converged with Reviewer 1 (R1.8) and the editor's specific request. Please see the response to R1.8 for the new HumanEval (code generation, 30 problems) and GSM8K (math reasoning, 30 problems) experiments, both demonstrating that the reproducibility gap widens precisely on tasks where small textual differences alter functionality or numerical correctness (Claude Sonnet 4.5 EMR = 0.393 on HumanEval and 0.063 on GSM8K). The non-AI/ML domain is also covered by the 10-abstract PubMed PM\textsubscript{2.5} probe (R3.4). We were not able to include open-ended \emph{creative-writing} tasks within the revision budget; we explicitly acknowledge this in the revised Limitations as a direction for future work, and note that the three new domains we did add (code, math, out-of-AI/ML extraction) are precisely the high-stakes regimes the editor highlighted in the major-revision letter.} \done

%--------------------------------------------------------------
\subsection*{R3.2 --- Quasi-isolation cannot be performed for proprietary APIs}
%--------------------------------------------------------------

\begin{revquote}
``Cloud deployment does not preclude reproducibility'' was tested with LLaMA 3 model via Together AI's cloud endpoint. Strictly speaking, this proof is not complete without verifying with ``Smarter'' APIs (GPT-4, Claude, Gemini). Since they are closed-source and proprietary, they cannot be locally deployed to compare against the cloud version. The infrastructure required to run them (like massive GPU clusters and speculative fast-paths) may currently make determinism an ``engineering impossibility'' at scale.
\end{revquote}

\response{The reviewer is correct that the quasi-isolation argument cannot, in principle, be extended to closed-source proprietary stacks: doing so would require local execution of those weights, which is not possible. We have appended a paragraph to the existing Limitations block in the Discussion (around line~410): ``The probe demonstrates that cloud deployment is not a \emph{sufficient} cause of LLM non-determinism\ldots\ We do not claim, and the result does not warrant claiming, that the proprietary stacks behind GPT-4, Claude, and Gemini could achieve the same level of reproducibility under their current production-infrastructure design constraints; their stacks are closed and we cannot run the symmetric `same-weights, two-stacks' comparison against them\ldots\ The Together~AI result is therefore offered as evidence for the weaker claim and as a partial existence proof that closed-source providers retain the option to publish reproducibility-prioritised execution modes.'' The Together~AI result remains sufficient evidence for the weaker, defensible claim while the stronger universal claim is appropriately qualified.} \done

\changes{Discussion §Limitations (around line 410): paragraph appended reformulating the quasi-isolation claim from ``cloud deployment does not preclude reproducibility'' (strong) to ``cloud deployment is not a sufficient cause of non-determinism'' (defensible), explicitly acknowledging the missing two-stack comparison for proprietary providers.}

%--------------------------------------------------------------
\subsection*{R3.3 --- Multi-turn coverage extension to GPT-4 and DeepSeek}
%--------------------------------------------------------------

\begin{revquote}
While five major API providers are tested, only two API providers (Claude and Gemini) are included in the more complex multi-turn and RAG experiments, which showed the most dramatic reproducibility failures. It is suggested to extend the multi-turn and RAG evaluations to OpenAI (GPT-4) and DeepSeek to confirm if the ``near-zero'' reproducibility found in Claude and Gemini is truly universal across all major cloud-served interactive pipelines.
\end{revquote}

\response{We agree this is a valuable extension. We have run the three-turn refinement protocol on the gpt-4o and DeepSeek Chat stacks (10 abstracts $\times$ 5 reps, 50 runs each):

\begin{itemize}[leftmargin=*]
  \item \textbf{gpt-4o multi-turn EMR = 0.09 [0.02, 0.16]} (10 abstracts, n=10 per-abstract EMR values).
  \item \textbf{deepseek-chat multi-turn EMR = 0.35 [0.13, 0.60]}.
\end{itemize}

This confirms the near-zero EMR pattern previously reported for Claude (0.04) and Gemini (0.01) is universal across major cloud-served interactive pipelines, not a Claude/Gemini-specific artefact. gpt-4o approaches the Gemini regime; deepseek-chat shows higher reproducibility but still drops more than four-fold from its single-turn level (0.84 $\rightarrow$ 0.35), consistent with error compounding across dialogue turns. We extended multi-turn coverage but did not extend the RAG-extraction protocol to gpt-4o and deepseek-chat within the revision budget; we note this scope explicitly in the Limitations section, and emphasise that single-turn extraction EMRs reported for both stacks (0.42 and 0.66 respectively) already place them in the regime where RAG would only amplify the gap, consistent with the Claude and Gemini RAG results in the original analysis. (Single-turn EMRs are tabulated in the revised manuscript Methods Table~1, ``Exact match rate (EMR) with 95\% bootstrap CI per deployment stack on the revision tasks''.)} \done

\changes{Revised Figure~2 now shows four API stacks for multi-turn (Claude, Gemini, gpt-4o, deepseek-chat) plus the three local stacks; updated Results §``Complex workflows amplify the reproducibility gap'' with the gpt-4o and deepseek-chat figures (above) and the four-fold drop observation. Data: \texttt{outputs/revision/runs/rev\_\{gpt-4o,deepseek-chat\}\_multiturn\_refinement\_*.json}.}

%--------------------------------------------------------------
\subsection*{R3.4 --- Limited domain (30 AI/ML abstracts)}
%--------------------------------------------------------------

\begin{revquote}
The study relies on 30 AI/ML abstracts. While statistically well-powered for the large effects observed, this limited domain may not reflect how non-determinism manifests in non-English texts or less structured fields.
\end{revquote}

\response{We thank the reviewer for raising this. We added 10 PubMed PM\textsubscript{2.5}/respiratory-health abstracts (single-turn structured extraction, C1, 5 reps $\times$ 8 stacks). Results confirm the local-API gap persists outside the AI/ML domain:

\begin{itemize}[leftmargin=*]
  \item Local + Together AI: EMR $\in$ [0.96, 1.00] (three stacks at 1.000).
  \item API stacks: deepseek-chat = 0.66, gemini-2.5-pro = 0.49, gpt-4o = 0.42, and \textbf{Claude Sonnet~4.5 = 0.010 [0.00, 0.03]}.
\end{itemize}

Claude's near-zero EMR on health abstracts mirrors its 0.020 EMR on the original AI/ML summarisation task --- the pattern is domain-transferable, not corpus-specific. This complements the larger 500-abstract analysis in our companion paper (Rover \& Tadano, 2026, submitted; cited but not duplicated here per R3.6). We acknowledge that non-English coverage remains a limitation and have stated this explicitly; we discuss it in the revised Limitations section as a direction for future work.} \done

\changes{New paragraph in Results §``Non-determinism varies substantially across providers'' reporting the PubMed PM\textsubscript{2.5} EMRs (above); new Extended Data Table row; explicit cross-reference to the companion paper for the full 500-abstract analysis.}

%--------------------------------------------------------------
\subsection*{R3.5 --- Mechanism mapping per stack}
%--------------------------------------------------------------

\begin{revquote}
While the paper lists six potential mechanisms for non-determinism, the manuscript would benefit from a more explicit discussion of which specific mechanisms are likely active in the Together~AI ``quasi-isolation'' case versus the OpenAI/Anthropic cases.
\end{revquote}

\response{We have added a new Table (revised manuscript Methods Table~2, label \texttt{tab:mechanisms}, ``Mapping of non-determinism mechanisms to deployment stacks''; rendered as \texttt{table*} for full-page width inside Methods §``Sources of non-determinism in distributed inference'') mapping each of the six mechanisms to each of five stacks (L-LLaMA3, C-LLaMA3 Together, API-GPT4, API-Claude, API-Gemini). Cells are labelled with one of: \checkmark (documented), ``likely'' (consistent with public information), ``possible'' (compatible with public information but not confirmed), or --- (precluded by stack design). A footnote states explicitly that attribution is inferential for closed-source providers. The accompanying Discussion paragraph uses this mapping to argue that the Together~AI residual variation is attributable predominantly to mechanisms 1 and~2 (FP non-associativity, BF16/FP16 mixed-precision), while GPT-4 / Claude / Gemini variation additionally implicates mechanisms 3--6 (tensor parallelism, FlashAttention, dynamic batching, speculative decoding).} \done

\changes{New \textbf{Methods Table~2} (label \texttt{tab:mechanisms}) inside §``Sources of non-determinism in distributed inference'' (\texttt{table*} float, full-width). Five columns × six rows; explicit footnote on inferential attribution for closed-source. Discussion paragraph using the mapping.}

%--------------------------------------------------------------
\subsection*{R3.6 --- Validation: are the 23 PM2.5 effects truly contradictory?}
%--------------------------------------------------------------

\begin{revquote}
The use of specific model snapshots (e.g., gpt-4-0613) means the exact EMR values may be ``point-in-time'' estimates that change as providers update their underlying infrastructure. Since the authors found high BERTScore F1 $>0.97$ despite low EMR, they should include a human evaluation or LLM-as-a-judge step to confirm if the 23 ``disappearing'' effect estimates in the PM\textsubscript{2.5} study are truly contradictory or just semantically varied.
\end{revquote}

\response{We thank the Reviewer for this important request and provide validation through a complementary triangulation strategy that preserves the integrity of our companion paper currently submitted:

\begin{itemize}[leftmargin=*]
  \item \textbf{Companion-paper validation (cited, not duplicated).} The 500-abstract analysis underlying the 23-effect finding is presented in detail in our companion paper (Rover \& Tadano, 2026, submitted to \emph{Research Synthesis Methods}), which independently performs systematic pairwise reproducibility analysis, silver-standard validation against an independent reasoning model (DeepSeek-R1), and meta-analytic propagation across 6 LLMs $\times$ 10 runs. To preserve that paper's contributions and avoid duplicate publication, the present manuscript does not reproduce its detailed tables; it cites them and reports only the aggregate propagation finding (23 study-level effect estimates).
  \item \textbf{Independent in-paper triangulation with two judges.} For a self-contained validation within the present manuscript, we ran a blinded two-judge LLM-as-judge protocol on a \emph{new} sample of \emph{30} cases drawn from the same corpus but \emph{distinct} from the 23 effects analysed in the companion paper. The two judges---Claude Opus~4.7 and gpt-4o---scored each pair independently against three pre-registered criteria: (a)~same direction of the reported effect (positive/null/negative); (b)~same magnitude within $\pm 20\%$; (c)~same 95\,\% confidence interval coverage. \textbf{Verdict counts}: Claude Opus 4.7 returned 22 \emph{truly contradictory} (Wilson 95\% CI 56--86\%), 3 \emph{semantically equivalent}, and 5 \emph{ambiguous}. gpt-4o returned 27 \emph{truly contradictory} (74--97\%), 3 \emph{semantically equivalent}, and 0 \emph{ambiguous}. \textbf{Inter-judge agreement: 76.7\,\% with Cohen's $\kappa$~=~0.29 (fair agreement, Landis--Koch).} Both judges classify the majority as substantive contradiction; the disagreement is concentrated in the ``ambiguous'' category where Claude is more conservative than gpt-4o. The lower bound on the truly-contradictory proportion (56\% Claude; 74\% gpt-4o) excludes a pure-cosmetic interpretation under BERTScore~F1~$>$~0.97. Total cost: USD~1.30 (Claude Opus on Anthropic + gpt-4o on OpenAI).
  \item \textbf{Point-in-time caveat addressed.} We also acknowledge that exact EMR values can drift as providers update infrastructure, and we note this in the revised Limitations. The gpt-4o snapshot used here (gpt-4o-2024-11-20 served via OpenAI Chat Completions) confirms the qualitative pattern persists at a current snapshot: gpt-4o single-turn extraction EMR = 0.42 on PubMed, 0.84 on HumanEval, 0.27 on GSM8K --- all dramatically below the corresponding local-stack performance.
\end{itemize}

This combined protocol provides validation in three layers: (i)~the large-corpus systematic analysis in our companion paper (cited only); (ii)~the independent in-paper two-judge LLM-as-judge triangulation (Claude Opus 4.7 + gpt-4o, N=30, inter-judge $\kappa$=0.29) that directly satisfies the ``human evaluation \emph{or} LLM-as-a-judge'' requirement of this comment; and (iii)~a current-snapshot drift check via gpt-4o-2024-11-20.

We note explicitly that we did \emph{not} run a human-rater validation for the present resubmission. The Reviewer's request offered ``human evaluation \emph{or} LLM-as-a-judge'' as a disjunction, and our two-judge LLM-as-judge protocol with cross-provider validation (Claude Opus 4.7 from Anthropic and gpt-4o from OpenAI---two distinct training families) addresses the ``LLM-on-LLM'' bias concern more directly than a single-judge design would. A dedicated methodological extension applying multi-rater human validation against a domain-expert silver standard to a larger sample is in preparation as a follow-up paper; an extension we will report in due course.} \done

\changes{New Results subsection ``Validation of the applied PM\textsubscript{2.5} divergence finding'' reporting the 30-case verdicts from both judges (Claude 22/3/5; gpt-4o 27/3/0; Cohen's $\kappa$=0.29); new Extended Data Table with per-case judge results (\texttt{outputs/revision/t3\_judge/}); gpt-4o snapshot drift comparison incorporated into the new EMR results. Detailed silver-standard, pairwise-disagreement, and meta-analytic tables remain in the companion paper to preserve its contributions. Multi-rater human validation is explicitly framed as future work in the revised Limitations.}

\bigskip\bigskip
\hrule
\bigskip

\noindent\emph{End of point-by-point response. Detailed change log accompanying this document: \texttt{response\_letter/02\_changes\_log.md}. Revised manuscript with track changes: \texttt{submission\_revision\_v1/ncomms\_main\_tracked.tex}. Clean revised manuscript: \texttt{article/ncomms\_main.tex}.}

\end{document}
